\section{Lompatan Bersyarat dan Percabangan}
\label{sec:jcc}

Lompatan bersyarat (\asm{Jcc}) memeriksa keadaan \eng{flag}: bila
syarat terpenuhi, eksekusi berpindah ke label tujuan; bila tidak,
eksekusi berlanjut ke instruksi berikutnya.

\subsection{Tabel Lompatan Bersyarat}

Lompatan bersyarat yang sering dipakai dirangkum pada
Tabel~\ref{tab:jcc}.

\begin{table}[H]
    \centering
    \caption{Lompatan bersyarat 8086 dan syaratnya.}
    \label{tab:jcc}
    \small
    \begin{tabular}{@{}p{3.2cm}p{4.2cm}p{5.8cm}@{}}
        \toprule
        \textbf{Instruksi} & \textbf{Syarat flag} &
        \textbf{Jenis bilangan} \\
        \midrule
        \asm{JE} / \asm{JZ} & \flag{ZF} $=1$ &
        Bertanda maupun tidak bertanda \\
        \addlinespace
        \asm{JNE} / \asm{JNZ} & \flag{ZF} $=0$ &
        Bertanda maupun tidak bertanda \\
        \addlinespace
        \asm{JG} / \asm{JNLE} & \flag{ZF} $=0$ dan
        \flag{SF} $=$ \flag{OF} & Bertanda \\
        \addlinespace
        \asm{JGE} / \asm{JNL} & \flag{SF} $=$ \flag{OF} &
        Bertanda \\
        \addlinespace
        \asm{JL} / \asm{JNGE} & \flag{SF} $\neq$ \flag{OF} &
        Bertanda \\
        \addlinespace
        \asm{JLE} / \asm{JNG} & \flag{ZF} $=1$ atau
        \flag{SF} $\neq$ \flag{OF} & Bertanda \\
        \addlinespace
        \asm{JA} / \asm{JNBE} & \flag{CF} $=0$ dan
        \flag{ZF} $=0$ & Tidak bertanda \\
        \addlinespace
        \asm{JAE} / \asm{JNB} & \flag{CF} $=0$ &
        Tidak bertanda \\
        \addlinespace
        \asm{JB} / \asm{JNAE} & \flag{CF} $=1$ &
        Tidak bertanda \\
        \addlinespace
        \asm{JBE} / \asm{JNA} & \flag{CF} $=1$ atau
        \flag{ZF} $=1$ & Tidak bertanda \\
        \bottomrule
    \end{tabular}
\end{table}

\subsection{Bertanda versus Tidak Bertanda}

Pemilihan \asm{Jcc} bergantung pada jenis bilangan. Setelah
\lstasm|CMP AX, BX|, perbandingan dapat ditafsirkan sebagai bilangan
bertanda atau tidak bertanda.

\begin{table}[H]
    \centering
    \caption{Pemetaan perbandingan ke lompatan yang benar.}
    \label{tab:jcc-signed}
    \small
    \begin{tabular}{@{}p{4.0cm}p{4.4cm}p{4.8cm}@{}}
        \toprule
        \textbf{Perbandingan} & \textbf{Bertanda} &
        \textbf{Tidak bertanda} \\
        \midrule
        \eng{dst} $>$ \eng{src} & \asm{JG} & \asm{JA} \\
        \addlinespace
        \eng{dst} $\geq$ \eng{src} & \asm{JGE} & \asm{JAE} \\
        \addlinespace
        \eng{dst} $<$ \eng{src} & \asm{JL} & \asm{JB} \\
        \addlinespace
        \eng{dst} $\leq$ \eng{src} & \asm{JLE} & \asm{JBE} \\
        \addlinespace
        \eng{dst} $=$ \eng{src} & \asm{JE}/\asm{JZ} &
        \asm{JE}/\asm{JZ} \\
        \addlinespace
        \eng{dst} $\neq$ \eng{src} & \asm{JNE}/\asm{JNZ} &
        \asm{JNE}/\asm{JNZ} \\
        \bottomrule
    \end{tabular}
\end{table}

Contoh, \hex{FF} dapat dibaca sebagai 255 (tidak bertanda) atau $-1$
(bertanda). Bila dibandingkan dengan 1, hasilnya berlawanan:
sebagai bilangan tidak bertanda \hex{FF} lebih besar, sedangkan
sebagai bilangan bertanda \hex{FF} lebih kecil.

\subsection{Percabangan if--else}

Pola \code{if--else} memerlukan satu lompatan untuk melewati blok
\code{else} setelah blok \code{if} selesai.

\begin{lstlisting}[language=MASM, numbers=none]
; if (AX == BX) { hasil = 1 } else { hasil = 0 }
CMP AX, BX
JNE BLOK_ELSE
MOV hasil, 1
JMP SELESAI
BLOK_ELSE:
MOV hasil, 0
SELESAI:
\end{lstlisting}

Alur tersebut ditunjukkan pada Gambar~\ref{fig:alur-if-else}.

\begin{figure}[H]
    \centering
    \begin{tikzpicture}[
        blok/.style={
            draw=obeBlue, line width=0.7pt, fill=obeBlueLight,
            rounded corners=2pt, minimum width=3.2cm,
            minimum height=0.8cm, align=center, font=\small,
            inner sep=3pt
        },
        uji/.style={
            draw=obeTeal, line width=0.7pt, fill=obeTealLight,
            diamond, aspect=2.4, minimum width=3.6cm,
            align=center, font=\scriptsize, inner sep=1pt
        },
        panah/.style={-{Stealth[length=2.2mm]}, line width=0.8pt,
            obeGray},
        ket/.style={font=\scriptsize, text=obeGray}
    ]
        \node[uji] (cmp) {AX $=$ BX?};
        \node[blok, below left=0.9cm and 1.6cm of cmp] (ya)
            {blok \code{if}\\hasil $=1$};
        \node[blok, below right=0.9cm and 1.6cm of cmp] (tidak)
            {blok \code{else}\\hasil $=0$};
        \node[blok, below=2.2cm of cmp] (selesai) {lanjut};

        \draw[panah] (cmp.west) -| node[ket, above left] {ya}
            (ya.north);
        \draw[panah] (cmp.east) -| node[ket, above right] {tidak}
            (tidak.north);
        \draw[panah] (ya.south) |- (selesai.west);
        \draw[panah] (tidak.south) |- (selesai.east);
    \end{tikzpicture}
    \caption{Alur percabangan \code{if--else} dengan \asm{CMP} dan
    \asm{JNE}.}
    \label{fig:alur-if-else}
\end{figure}

\begin{peringatanbox}
\vspace{1.1em}
Jangan menyisipkan instruksi yang mengubah \eng{flag} di antara
\asm{CMP} dan \asm{Jcc}. Instruksi seperti \asm{INC}, \asm{DEC},
\asm{ADD}, atau \asm{LOOP} dapat merusak syarat lompatan. Bila perlu
menyimpan nilai, gunakan \asm{MOV} yang tidak mengubah \eng{flag}.
\end{peringatanbox}
